\section{Primera iteración de la Fase 6: Mega y ESP32, comparación y estrategia}
\label{sec:F6I01}

Al disponer de ambos sistemas para el desarrollo en el sistema externo, en esta iteración comparo las dos plataformas y sus posibilidades para llevar el motor al hardware de ambas máquinas: un Arduino Mega 2560 con pantalla usando un shield con todos los pines necesarios, y un ESP32 (que es el destino real del proyecto) con pantalla SPI, y verifica si es posible usar ambas arquitecturas como destino del proyecto.

\subsection{Las dos plataformas}
\label{sec:F6I01Plataformas}

\begin{table}[H]
	\centering
	\begin{tabular}{lp{5.3cm}p{5.3cm}}
		\toprule
		 & Arduino Mega 2560 & ESP32 DevKitC V2 \\
		\midrule
		Procesador & AVR de 8 bits, 16\,MHz & Xtensa de 32 bits, 2 núcleos, hasta 240\,MHz \\
		Flash & 256\,KB & 4\,MB (en la placa de pruebas) \\
		RAM & 8\,KB & unos 520\,KB, de los que quedan unos 300\,KB utilizables \\
		PSRAM & No & No \\
		Lógica & 5\,V & 3,3\,V \\
		Pantalla & Shield TFT de 2,4\,pulgadas de MCUFRIEND con lector de SD, bus paralelo de 8 bits, 240$\times$320 QVGA, color de 16 bits & ARCELI ILI9341 de 2,8\,pulgadas, SPI, con lector de SD \\
		Datos & En \texttt{PROGMEM}, se leen con \texttt{pgm\_read\_*}; por encima de 64\,KB hacen falta punteros <<far>> & Flash mapeada en memoria; se leen como memoria normal \\
		Pantalla: envío & Escritura directa a puertos o biblioteca; sin DMA & SPI con DMA y búferes \\
		Tiempo y sistema & \texttt{millis()}, \texttt{micros()}, \texttt{loop()} & \texttt{esp\_timer}, FreeRTOS, \texttt{loop()} de Arduino \\
		\bottomrule
	\end{tabular}
	\caption{Comparación de las plataformas}
	\label{tab:F6I01Plataformas}
\end{table}

\begin{table}[H]
	\centering
	\begin{tabular}{lp{5.3cm}p{5.3cm}}
		\toprule
		& Arduino Mega 2560 & ESP32 DevKitC V2 \\
		\midrule
		Cartucho & Compilado dentro del programa (array en \texttt{PROGMEM}) & Partición de datos de la flash \\
		Fondo & No es posible (no cabe en la RAM ni se puede mapear cómodamente) & Sí, desde la flash \\
		Carga de cartuchos & Posibilidad de usar la SD & Posibilidad de usar la SD,  Bluetooth o Wifi\\
		\bottomrule
	\end{tabular}
\caption{Donde irán los datos en cada sistema}
\label{tab:F6I01PlataformasDatos}
\end{table}


\subsection{Qué se puede hacer en cada una}
\label{sec:F6I01Posible}

\begin{itemize}
	\item \textbf{Mega, posible sin incluir el fondo:} el motor completo \textbf{parece posible} para un juego de la talla del Parachute (estimación, solo se puede confirmar ejecutando el desarrollo):
	\begin{itemize}
		\item \textbf{Flash:} el GWC del Parachute son unos 7\,KB y el código del motor con las bibliotecas rondará 20 a 40\,KB, de 256\,KB.
		\item \textbf{RAM:} el estado de ejecución (marcas de visible y cambiar de unos 54 segmentos, rutas activas, contadores, jugador, temporizadores y posición del generador aleatorio) ronda 1\,KB; con la biblioteca de pantalla, la pila y el puerto serie, debería quedar por debajo de 3\,KB de los 8\,KB.
		\item \textbf{Velocidad:} la lógica sobra; lo costoso es pintar, unos pocos milisegundos por sprite, y en cada iteración solo cambian unos pocos.
		\item \textbf{Simplicidad}: Al no incluir el fondo no hay que pensar en recortes, el motor gráfico maneja los sprites de forma que se borran y escriben directamente en la pantalla sin preocuparse de conflictos con fondo o entre sí, las imágenes son una rsitra de ceros y unos, para dibujarlos donde hay un 0 no se hace nada, donde hay un 1 se pinta un punto negro, para eliminarlos si hay un cero no se hace nada, si hay un 1 se pinta un punto del color del fondo elegido.
	\end{itemize}
	\item \textbf{Mega, lo que no se puede:} imagen de fondo, cartuchos intercambiables (la memoria de programa no se escribe con facilidad en ejecución) y el selector de la SD; juegos con muchos más segmentos o rutas toparían con la RAM.
	\item \textbf{ESP32:} todo lo anterior y además el fondo, el selector, y mucho margen de RAM y de velocidad. Es el destino final.
\end{itemize}

\subsection{Estrategia}
\label{sec:F6I01Estrategia}

\begin{itemize}
	\item \textbf{Un solo código con compilación condicional.} El programa se escribe una sola vez en C/C++ con el entorno de Arduino (el mismo proyecto, p.~ej. con PlatformIO) y se compila expresamente para cada aparato. Las diferencias se aíslan con \texttt{\#if defined(ARDUINO)} y \texttt{\#if defined(ESP32)} en una pequeña capa de plataforma.
	\item \textbf{Lo común}: casi todo, el director, el motor de estados, el lector del GWC y la lógica del juego.
	\item \textbf{Lo que cambia}: la capa de plataforma y el motor gráfico (aunque solo la parte que muestra y borra los sprites, el núcleo base es común)
\end{itemize}

\begin{table}[H]
	\centering
	\begin{tabular}{lp{3.2cm}p{3.6cm}p{3.6cm}}
		\toprule
		Pieza & Común & Mega & ESP32 \\
		\midrule
		\texttt{juegoDirector} & Bucle y reloj & \texttt{millis()}, \texttt{micros()} & \texttt{esp\_timer} \\
		\texttt{juegoMotorEstados} & Todo & -- & -- \\
		\texttt{baseCartucho} (lector) & Estructura y validación (cabecera, CRC) & Lee de \texttt{PROGMEM} & Lee de la flash mapeada \\
		\texttt{basePantalla} & Funciones públicas & Puertos o biblioteca & SPI con DMA \\
		\texttt{juegoMotorGrafico} & Interfaz (\texttt{mgIniciar}, \texttt{mgFijarSegmento}, \texttt{mgPresentar}) & Fondo liso: se pinta la máscara & Fondo con imagen: compositor) \\
		\texttt{configHardware} & -- & Pines, resolución & Pines, resolución, reloj SPI \\
		\bottomrule
	\end{tabular}
	\caption{Qué es común y qué cambia en cada plataforma}
	\label{tab:F6I01Comun}
\end{table}

\begin{itemize}
	\item \textbf{Acceso a los datos del cartucho siempre por funciones} (\texttt{cartLeer8}, \texttt{cartLeer16}, \texttt{cartLeer32}), nunca por punteros ni estructuras en RAM. En el ESP32 son una lectura de memoria; en el Mega, un \texttt{pgm\_read\_*}. El lector y los motores son idénticos en las dos. Es la decisión que más trabajo ahorra y no cuesta nada si se toma antes de escribir código.
	\item \textbf{El motor gráfico tiene la misma interfaz en las dos plataformas;} solo cambia su interior (ver el siguiente apartado).
	\item \textbf{Reglas para la RAM del Mega que se aplican a todo el código:} nada de \texttt{String} ni \texttt{malloc}; arrays de tamaño fijo con constantes máximas; si hay textos de depuración (tema a evitar por consumo de memoria) se usará \texttt{F("...")}; tipos de ancho fijo (\texttt{uint8\_t}, \texttt{uint16\_t}, \texttt{uint32\_t}), porque el \texttt{int} del AVR es de 16 bits. Cumplirlas no perjudica al ESP32.
\end{itemize}

\subsubsection{El motor gráfico en las dos etapas}

\begin{itemize}
	\item \textbf{Fondo liso (Mega):} se pinta el fondo con \texttt{ColorFondo} (Bloque~0, F5-I13) y los sprites son máscaras negras. Para presentar un segmento se pintan en negro los píxeles de su máscara; para borrarlo, esos mismos píxeles con el color de fondo. Como los píxeles de dos segmentos nunca coinciden, borrar uno no daña a otro: no hay composición, ni lista de vecinos, ni búfer de RAM. No hace falta XOR (que además exigiría leer la memoria de vídeo, algo que estas pantallas no permiten).
	\item \textbf{Fondo con imagen (ESP32):} el fondo vive en la flash y se compone cada rectángulo sucio con el fondo y los sprites visibles que lo cortan.
	\item Los dos conservan la misma interfaz hacia el director y el motor de estados, de modo que pasar del Mega al ESP32 solo cambia el interior del motor gráfico y la capa de plataforma.
\end{itemize}

\subsection{Hoja de ruta de la Fase 6}
\label{sec:F6I01Ruta}

\begin{enumerate}
	\item \textbf{Mega sin fondo.}
	\begin{enumerate}
		\item \textbf{MP1} Prueba gráfica mínima: el paracaidista cae \underline{sobre una pantalla sin fondo}, con el cartucho real del Parachute, el director y un motor de estados mínimo.
		\item \textbf{MP2} Motor de estados completo (teclado, rutas, aleatorios, colisiones, contadores y fin del juego, las Fases 2 y 3) llevado al Mega, con botones físicos.
	\end{enumerate}
	\item \textbf{ESP32 con fondo.} 
	\begin{enumerate}
		\item \textbf{EP1} Prueba gráfica mínima: el paracaidista cae \underline{sobre una pantalla con un fondo}, con el cartucho real del Parachute, el director y un motor de estados mínimo.
		\item \textbf{EP2} Desarrollo completo: pantalla y reloj SPI, el lector del cartucho sobre la flash mapeada, el fondo y el compositor, con la prueba del segmento móvil. El motor de estados no se toca: es el mismo de la etapa 1.
	\end{enumerate}
	\item \textbf{Tercera Etapa (se desarrolla en la Fase 7):} selector de cartuchos desde la SD.
\end{enumerate}

El orden adecuado de desarrollor será MP1, EP1, MP2, EP2, ya que se comparte en ambos casos la mayor parte del desarrollo, los motores serán prácticamente iguales en ambos aparatos, lo que cambia es la parte gráfica.

\subsection{Equipo necesario para la etapa 1}
\label{sec:F6I01Equipo}

\begin{itemize}
	\item \textbf{Arduino Mega 2560 R3 con el shield TFT de 2,4\,pulgadas de MCUFRIEND con SD.} Es un shield de formato UNO con bus paralelo de 8 bits. 
	\begin{itemize}
		\item \textbf{Resolución:} 320$\times$240, que es la mínima que admite el GWC.
		\item \textbf{Controlador:} los shields de MCUFRIEND montan varios dodelos de controlador diferente, no visible en la placa (ILI9341, ILI9325, ST7781, \ldots). La biblioteca \texttt{MCUFRIEND\_kbv} lo identifica leyendo su código.
		\item \textbf{Pines que ocupa:} los de datos y control en los conectores del formato UNO (digitales 2 a 9 para los datos y analógicos A0 a A4 para el control), y los de la SD en los digitales 10 a 13, que en el Mega \textbf{no coinciden} con su SPI físico (50 a 53). 
		\item \textbf{Pines libres:} los analógicos 8 a 15 y los digitales 22 a 49 quedan libres, los usados para los botones serán de los digitales por comodidad.
		\item \textbf{Biblioteca:} \texttt{MCUFRIEND\_kbv} con \texttt{Adafruit\_GFX}; según su documentación incluye escritura directa a los puertos del Mega, que es lo que da velocidad (a comprobar midiendo).
	\end{itemize}
	\item \textbf{Botones:} de 2 a 5 botones de juego más los de sistema (Start, Stop, Reloj; F5-I10), conectados a los pines libres 22 a 49. Con las resistencias de arrastre internas (\texttt{INPUT\_PULLUP}) no hacen falta resistencias externas: el botón conecta el pin a masa. Si se usan resistencias externas, la clásica de 10\,k$\Omega$ de arrastre a positivo con el botón a masa funciona igual y se lee con la misma lógica invertida.
	\item \textbf{Anti-rebote} por programa en el director.
\end{itemize}

\subsection{Riesgos y puntos abiertos}
\label{sec:F6I01Abiertos}

\begin{itemize}
	\item \textbf{RAM del Mega:} es el riesgo principal. Se mide desde el primer día (memoria libre al arrancar y mínima durante el juego) y las reglas del apartado de estrategia se cumplen siempre.
	\item \textbf{Velocidad de dibujo:} depende de la pantalla y de la biblioteca. Si la de serie fuera lenta, se sustituye por escritura directa a puertos solo en \texttt{basePantalla}.
	\item \textbf{Modelo de la pantalla del Mega y biblioteca:} controlador, resolución, orientación y pines. Se comprueba en el paso 2.
	\item \textbf{No optimizar de más para el Mega:} que el código quepa en 8\,KB es una ventaja, pero el destino es el ESP32. La única exigencia propia del Mega es no usar memoria dinámica.
	\item \textbf{Entorno:} Arduino IDE o PlatformIO; para dos plataformas con un mismo código es más cómodo PlatformIO (un entorno por plataforma). Se decide tras pruebas antes de comenzar.
	\item \textbf{Cartucho en el Mega:} se vuelca el GWC a un fichero de cabecera (\texttt{.h}) con un guion (tipo \texttt{xxd -i}); a medio plazo, un botón del Designer podría generarlo.
\end{itemize}

\subsection{Criterios de aceptación de la decisión}
\label{sec:F6I01Criterios}

\begin{itemize}
	\item La Fase~6 se hace en etapas: Mega sin fondo y ESP32 con fondo, primero la prueba simple y luego la completa.
	\item Se definirán completamente la capa de plataforma y la interfaz común del motor gráfico, de modo que el trabajo se reutiliza íntegro para cada máquina
	\item \textbf{Importante}: Si las pruebas simples de movimiento del paracaidista no se ven correctas, se abandona el intento de usar el Mega y se pasa al ESP32 directamente.
\end{itemize}

\clearpage
